\problem{P35}{剩余区间中的次序位置}{条件缩放}
从 $(0,1)$ 独立均匀抽取 $n$ 个样本，次序统计量记为 $U_{(1)}<\cdots<U_{(n)}$。对 $1\le i<j\le n$，令
\[
W=\frac{U_{(j)}-U_{(i)}}{1-U_{(i)}}.
\]
证明 $W$ 与 $U_{(i)}$ 独立，求 $W$ 的分布、均值和方差，并计算 $\Prob(U_{(j)}>(1+U_{(i)})/2)$。
\solution 给定 $U_{(i)}=u$，上方 $n-i$ 个样本的无序集合服从 $(u,1)$ 上独立均匀样本的规律。这可由次序统计量的联合密度为常数，并在 $u$ 两侧分别分解积分区域得到。将上方样本缩放为 $(v-u)/(1-u)$ 后，便是 $m=n-i$ 个标准均匀样本。

$W$ 为其中第 $r=j-i$ 个次序统计量。通过在 $w$ 下方放 $r-1$ 个样本、在 $w$ 附近放一个样本、其余放在上方，得到
\[
f_W(w)=\frac{m!}{(r-1)!(m-r)!}
 w^{r-1}(1-w)^{m-r},\quad0<w<1.
\]
这个条件密度不依赖 $u$，所以 $W$ 与 $U_{(i)}$ 独立，且 $W\sim\operatorname{Beta}(r,m-r+1)$。用贝塔积分计算一、二阶矩，得到
\[
\boxed{\E W=\frac r{m+1},\qquad
\Var W=\frac{r(m-r+1)}{(m+1)^2(m+2)}.}
\]
$W>1/2$ 等价于 $m$ 个缩放样本中至多 $r-1$ 个落在左半区间，因此
\[
\boxed{\Prob\!\left(U_{(j)}>\frac{1+U_{(i)}}2\right)
=2^{-m}\sum_{k=0}^{r-1}\binom mk.}
\]
\insight 绝对位置互相依赖，条件缩放后的相对位置却可以独立。找出条件密度中能被消去的尺度，比直接算协方差更能揭示结构。
